\documentclass[10pt]{article}
\usepackage[margin=0.7in]{geometry}
\usepackage{amsmath,amssymb,booktabs,graphicx}
\usepackage[T1]{fontenc}
\graphicspath{{figures/}}
\title{RISE: Supplementary Experimental Analyses}
\author{Research supplement to the supplied manuscript}
\date{23 September 2026}
\begin{document}
\maketitle
\noindent This supplement reports completed exploratory analyses and matched controls.
RISE is named SISA in the implementation; AM and PA correspond to decay and
rotation. Native A denotes an independently trained backbone, whereas B retains
auxiliary content projections. C, D, and E denote AM-only, PA-only, and Full RISE.
The original 16-cell, seed-20262027 main table is unchanged. New controls use the
matched seed-20262028 protocol, not replacements selected for the original table.
Intervals quantify test-user sampling uncertainty conditional on fitted models,
not training-seed variation. The accompanying Chinese dossier contains the full
execution ledger, action-specific subgroup support, rejected diagnostics and
reproduction records. These materials are proposed manuscript additions; the
supplied manuscript PDF has not been edited.
\section{Sample-conditional effects and statistical validation}
\input{paper/native_merrec_analysis.tex}
\clearpage
\input{paper/component_analysis.tex}
\clearpage
\section{Architectural controls, efficiency and mechanism diagnostics}
\input{paper/extended_analyses.tex}
\clearpage
\section{Additional native-control figures}
\begin{figure}[ht]
\centering
\includegraphics[width=.96\textwidth]{figures/native_merrec_effect_by_history.pdf}
\caption{Full RISE versus the native HyFormer backbone on MerRec by effective history length. Intervals condition on the trained models and use paired user-cluster influence functions.}
\end{figure}
\begin{figure}[ht]
\centering
\includegraphics[width=.96\textwidth]{figures/native_merrec_effect_by_frequency.pdf}
\caption{Native-controlled MerRec effects by training-target frequency. Frequency zero is not proof of strict cold start.}
\end{figure}
\end{document}
